\section{Περιλήψεις}
\label{Summaries}
\subsection{Summary in English}

This thesis aims to solve the problem of human speech recognition using visual cues, more specifically a video of the speech in question. This is a problem very similar to the classic problem of Automatic Speech Recognition (ASR) and is thereby often called Visual Speech Recognition (VSR) or Audio-Visual Speech Recognition (AVSR), in the cases where both audio as well as video cues are used to predict speech. The field of ``lip reading'' has grown immensely in the last few years, especially since the advent of Transformers and the ability to train large artificial intelligence models using deep learning. Recent studies tend to divide the problem of VSR into two separate sub-problems; first, the prediction of phonemes (the smallest units containing phonetic information for human speech) and second, the combination of the predicted phonemes so as to predict words and, ultimately, phrases that the original speaker uttered.

During this particular work, more emphasis is placed on solving the first sub-problem, more specifically the problem of phoneme prediction using the video of a speaker. An artificial intelligence model is developed, which takes in as input the video of the speaker, centered around the area of the lips and without auditory information and returns its predicted phonemes per video frame. This model contains a 3D Convolutional Neural Network (CNN) responsible for the extraction of visual features, as well as a BiLSTM architecture as a Sequence Encoder. Moreover, a similar model is implemented which takes in as input both the video of the area of the lips as well as audio cues. This audio-visual model is very similar to the visual model, since it also uses CNNs for the feature extraction process, for both video and audio. A cross-attention mechanism is then applied to these extracted features, followed by a fusion block, which is gated by the reliability of each modality at any given time. Lastly, a BiLSTM architecture is again used as the Sequence Encoder.

The experimental process of this work consists of three separate experiments. During the first, the visual model is benchmarked in a constrained environment, in terms of both the vocabulary spoken by the speaker as well as the syntax they employ. The first sub-experiment is an ablation study, which produces the final proposed method as the model with the best results. The second sub-experiment studies the performance of the proposed method more in-depth, also using an algorithm for word prediction based on the predicted phonemes, which takes advantage of the constraints present in the chosen dataset. The goal of the second experiment is to test the performance of the proposed method in an unconstrained environment, using an ``in the wild'' dataset. However, since access to a particular popular ``in the wild'' dataset was not granted, the completion of this experiment was not made possible and the performances of various state of the art models were included for the sake of completeness. The completion of this experiment is mentioned as the most important future expansion of this study. Lastly, the third experiment tested the performance of the proposed audio-visual method, in the same constrained environment, but this time including noise in both modalities. During the first sub-experiment, the results of increasing audio SNR are studied, whereas during the second sub-experiment, the model is tested during conditions where the visual noise is increased.

\newpage
\subsection{Περίληψη στα Ελληνικά}

Η συγκεκριμένη εργασία ασχολείται με το πρόβλημα της αναγνώρισης της ανθρώπινης ομιλίας αξιοποιώντας οπτικά μέσα, δηλαδή βίντεο της ομιλίας. Είναι πρόβλημα συγγενικό με το κλασικό πρόβλημα της αυτόματης αναγνώρισης ομιλίας (Automatic Speech Recognition - ASR) και αποκαλείται συχνά Οπτική Αναγνώριση Ομιλίας (Visual Speech Recognition - VSR) ή Οπτικοακουστική Αναγνώριση Ομιλίας (Audio-Visual Speech Recognition - AVSR), σε περίπτωση που χρησιμοποιούνται και ηχητικά στοιχεία, πέραν των οπτικών. Ο τομέας της ``αναγνώρισης χειλιών'' (lip reading) έχει γνωρίσει τον τελευταίο καιρό σημαντική ανάπτυξη, ιδιαίτερα με την έλευση των Transformers και τη δυνατότητα εκπαίδευσης μοντέλων αξιοποιώντας τη βαθιά μάθηση. Πρόσφατες μελέτες τείνουν να διαχωρίσουν το πρόβλημα της αναγνώρισης της ομιλίας σε 2 υπο-προβλήματα: αφενός την αναγνώριση των φωνημάτων (των ελάχιστων φωνητικών τμημάτων της ανθρώπινης ομιλίας) και αφετέρου τον συνδυασμό των προβλεπόμενων φωνημάτων για την πρόβλεψη λέξεων και, τελικά, φράσεων.

Στα πλαίσια της εργασίας μελετάται η επίλυση του πρώτου υπο-προβλήματος, δηλαδή της πρόβλεψης φωνημάτων από βίντεο ενός ομιλητή. Αναπτύσσεται μοντέλο τεχνητής νοημοσύνης το οποίο δέχεται ως είσοδο βουβό βίντεο της περιοχής των χειλιών και επιστρέφει ως έξοδο τις προβλέψεις του για το ποιο φώνημα εκφέρεται σε κάθε frame του βίντεο. Το μοντέλο αυτό αξιοποιεί ένα 3D συνελικτικό νευρωνικό δίκτυο (CNN) για την εξαγωγή features από το βίντεο και BiLSTM ως Sequence Encoder. Παράλληλα, σχεδιάζεται μοντέλο που δέχεται ως είσοδο τόσο το βίντεο της περιοχής των χειλιών όσο και ηχητικά στοιχεία. Το οπτικο-ακουστικό μοντέλο ακολουθεί λογική παρόμοια με αυτή του οπτικού μοντέλου, αξιοποιώντας CNNs για την εξαγωγή των οπτικών και ακουστικών χαρακτηριστικών. Στα χαρακτηριστικά εξαγόμενα από κάθε modality εφαρμόζεται μηχανισμός cross-attention, ενώ τα αποτελέσματα αυτού συνενώνονται με βάση την αξιοπιστία του κάθε modality για κάθε χρονική στιγμή. Τέλος, εφαρμόζεται και πάλι BiLSTM ως Sequence Encoder.

Η πειραματική διαδικασία της εργασίας αυτής αποτελείται από 3 πειράματα. Κατά το πρώτο πείραμα, δοκιμάζεται η επίδοση του οπτικού μοντέλου σε περιβάλλον περιορισμένο τόσο ως προς το λεξιλόγιο όσο και ως προς το συντακτικό του. Στο πρώτο υπο-πείραμα, πραγματοποιείται μία μελέτη αφαίρεσης, από την οποία προκύπτει η προτεινόμενη μεθοδολογία ως αυτή με τις βέλτιστες επιδόσεις. Στο δεύτερο υπο-πείραμα, μελετάται σε μεγαλύτερο βάθος η επίδοση του προτεινόμενου μοντέλου, αξιοποιώντας και αλγοριθμικό μηχανισμό πρόβλεψης λέξεων από τα φωνήματα. Στόχος του δεύτερου πειράματος είναι η μελέτη της επίδοσης της προτεινόμενης μεθοδολογίας σε περιβάλλον μη περιορισμένο από το λεξιλόγιο και το συντακτικό, δηλαδή σε in the wild συνθήκες. Καθώς δεν κατέστη δυνατή η απόκτηση πρόσβασης σε διαδεδομένο in the wild dataset του χώρου, σημειώθηκαν οι επιδόσεις των state of the art μοντέλων για λόγους πληρότητας και αναφέρθηκε η διεκπεραίωση του συγκεκριμένου πειράματος ως η σημαντικότερη μελλοντική επέκταση. Τέλος, στο 3ο πείραμα μελετήθηκε η επίδοση της προτεινόμενης οπτικο-ακουστικής μεθοδολογίας σε περιβάλλον περιορισμένο αλλά επηρεασμένο από θόρυβο, τόσο ηχητικό όσο και οπτικό. Στο πρώτο υπο-πείραμα μελετάται η επίδραση της αύξησης του SNR του ηχητικού σήματος στην επίδοση του μοντέλου, ενώ στο δεύτερο υπο-πείραμα παρουσιάζονται τα αποτελέσματα της αύξησης του οπτικού θορύβου στις προβλέψεις του μοντέλου.